\documentclass{article}
\usepackage[utf8]{inputenc}
\usepackage[T1]{fontenc}
\usepackage{geometry}
\geometry{
    a4paper,
    left=2.5cm,
    right=2.5cm,
    top=2.5cm,
    bottom=2.5cm
}
\usepackage{graphicx}
\usepackage{float}
\usepackage{hyperref}
\usepackage[english]{babel}
\usepackage{amsmath}
\usepackage{amssymb}
\usepackage{array}
\usepackage{booktabs}
\usepackage{caption}
\usepackage{subcaption}
\usepackage{enumitem}
\usepackage{fancyhdr}
\usepackage{titlesec}
\usepackage{url}
\usepackage{textcomp}
\usepackage{xcolor}
\usepackage{listings}
\usepackage{multirow}
\usepackage{longtable}
\usepackage{tabularx}
\usepackage{times}

% Set up page style
\pagestyle{fancy}
\fancyhf{}
\rfoot{\thepage}
\renewcommand{\headrulewidth}{0pt}

% Title formatting
\titleformat{\chapter}[display]{\normalfont\bfseries}{\chaptertitle\thechapter}{0pt}{\Huge}
\titlespacing*{\chapter}{0pt}{50pt}{40pt}

% Document info
\title{Template Document}
\author{Academic Support Agent}
\date{\today}

\begin{document}
\maketitle

%% ep grafe OPCIONAL -> si no se desea, comentar la siguiente linea
\epigrafe
%%%%   NO TOCAR LA SIGUIENTE LINEA !!!
\materialfijoinicial
% pr logo / preface OPCIONAL -> si no se desea, comentar la siguiente linea
\prefacio
% abreviaturas y s mbolos / glossary OPCIONAL -> si no se desea, comentar la siguiente linea
\glosario
%%%%   NO TOCAR LA SIGUIENTE LINEA !!!
\pagenumbering{arabic}%
%
%%%% Recordar que los objetivos deben presentarse expl citamente de forma obligatoria, como una secci n en la introducci n o como un cap tulo aparte !!!!
% Content from introduccion would be included here


%%
%%%% Si queremos un ep grafe antes del cap. 2, por ejemplo, descomentamos la siguiente linea
%\chapterepigraph{Je pense, donc je suis}{Discours de la m thode}{R.\ Descartes}
%%
% Content from cap2 would be included here


% Content from cap3 would be included here


%%%% A danse tantos cap tulos como uno necesite
%
%%%% Recordar que en el formato de tesis por compendio, debe haber un cap tulo de discusi n global del trabajo que se incluye en los art culos antes de las conclusiones  !!!!
%
%% Content from cap3 would be included here

%
%
\conclusionesperspectivas
%%% SI hay alg n ap ndice (sino quitar las tres lineas siguientes) 

% Content from apendice1 would be included here


%% Content from apendice2 would be included here

%
%%% terminan ap ndices

%%%% BIBLIOGRAF A (se puede cambiar el estilo ... no se da soporte en ese caso...)

\phantomsection
\addcontentsline{toc}{chapter}{\bibname}%
%% La bibliograf a est  en Plantilla.bib. Se recomienda usar un gestor de bibliograf a ( p.ej. JABREF )



%%%  NDICE ALFAB TICO
\ifthenelse{\equal{\noindicealfabetico}{\null}}{\phantomsection\addcontentsline{toc}{chapter}{\indexname}\input Plantilla.ind\relax}{\relax}%
%
%%% en casos especiales (solo a veces es necesario en espa ol) hay que usar esto.   Preguntar !!!
%\input Plantilla.inx
%
%
% RES MENES (obligatorios). Pueden ponerse despu s de \materialfijoinicial
%  Uno tiene que ser la traducci n del otro
%
\resumeneningles % M XIMO UNA PAGINA !! (el fichero es summary.tex)
\resumenenespanyol % M XIMO UNA PAGINA !! (el fichero es resumen.tex)
%
%%%% Y SE ACABO --   NO A ADIR NADA MAS !!!
%
\end{document}